\section{Research published in 2026}
The review includes nine papers first submitted to arXiv between February and September 2026. Submission dates were checked separately from revision dates; exact consulted versions appear in the references. Selection targeted coordination, context transfer, skill assignment and measurable efficiency. This is a focused technical review, not an exhaustive systematic review. Other studies' scores are not combined with the experiment below.

\subsection{Execution topology and authority}
Kulkarni and Kulkarni compare sequential, parallel, supervisor--worker and reflexive workflows for financial document extraction. Their descriptions motivate the four pattern families studied here. The domain and January 2025 model/pricing setup constrain transfer to a personal assistant; we do not reproduce their configurations or reported performance.\source{https://arxiv.org/html/2603.22651v1}{Kulkarni and Kulkarni, arXiv:2603.22651v1, 2026}

DESBench examines centralized, hierarchical, heterarchical and holonic coordination in industrial scheduling. Its implemented heterarchy uses mediated bidding. It expands the design question from graph shape to who may allocate work, reject commitments and resolve conflicts; its industrial simulator is a different evaluation environment.\source{https://arxiv.org/html/2605.13172v1}{Wang et al., arXiv:2605.13172v1, 2026}

ParaGUI dispatches concurrent GUI workers in rounds and integrates their summaries. It provides evidence on decomposable interactive workloads, while identifying worker execution failures as a bottleneck. Its visual-history controls and multi-desktop setting make it unsuitable as a direct numerical baseline for this text-only experiment.\source{https://arxiv.org/html/2607.22689v1}{Yu et al., arXiv:2607.22689v1, 2026}

\subsection{Information transfer and evaluation}
OrchBench represents dependencies and cross-agent transfers, then simulates quality, makespan and token cost. Its retention mechanism makes information loss a coordination variable. Our benchmark instead executes every worker against the API; simulated worker costs or outcomes are not substituted for observed usage or answers.\source{https://arxiv.org/html/2607.25656v1}{Ren et al., arXiv:2607.25656v1, 2026}

CRAFT gives agents complementary partial views of a construction problem. It distinguishes individual reasoning from effective communication and describes repeated unproductive corrections. It motivates preserving source evidence and assessing the final result. Because every worker here receives the original task, we do not test its partial-observation problem.\source{https://arxiv.org/html/2603.25268v2}{Nath, VanderHoeven and Krishnaswamy, arXiv:2603.25268v2, 2026}

\subsection{Skills and scoped workers}
\textit{Subagents vs Agent Skills} distinguishes procedural packages with explicit input/output interfaces from loosely structured skill text. Its main evaluation uses a selected 64-task subset and reports lower peak context with greater total token consumption. This supports testing payloads separately from topology.\source{https://arxiv.org/html/2609.09233v1}{Piriyakulkij et al., arXiv:2609.09233v1, 2026}

EvoAgent combines hierarchical delegation with keyword, embedding and model-based skill matching. Its domain-specific evaluation uses model judges; those scores do not establish exact execution reliability. We use its separation of retrieval and delegation as a reference, without implementing autonomous learning or reproducing its retrieval cascade.\source{https://arxiv.org/html/2604.20133v3}{A. Zhang et al., arXiv:2604.20133v3, 2026}

Swarm Skills packages roles, workflows, limits and dependencies as reusable coordination assets. Its qualitative case study leaves broad conformance testing unresolved. We treat portable workflow descriptions as a design proposal, not evidence that arbitrary multi-agent frameworks can interchange them without validation.\source{https://arxiv.org/html/2605.10052v2}{X. Zhang et al., arXiv:2605.10052v2, 2026}

Xu and Yan survey progressive disclosure: discover metadata, load chosen skill instructions, then retrieve supporting resources. This distinguishes a catalog from the content injected into each context. As a survey, it informs the loading policy but supplies no independent performance result for our architectures.\source{https://arxiv.org/html/2602.12430v4}{Xu and Yan, arXiv:2602.12430v4, 2026}
